\documentclass{article}
\usepackage{amsmath, amssymb, amsthm}
\usepackage{hyperref}
\usepackage{geometry}
\geometry{margin=1in}

\title{Erd\H{o}s Problem \#468}
\date{Last updated: 2025-08-31}
\author{Source: erdosproblems.com}

\begin{document}
\maketitle

\noindent\textbf{Status:} OPEN \quad \textbf{No prize}

\noindent\textbf{Tags:} number theory, divisors

\medskip

\noindent\textbf{Problem Statement:}

For any $n$ let $D_n$ be the set of sums of the shape $d_1,d_1+d_2,d_1+d_2+d_3,\ldots$ where $1&#60;d_1&#60;d_2&#60;\cdots$ are the divisors of $n$. What is the size of $D_n\backslash \cup_{m&#60;n}D_m$? If $f(N)$ is the minimal $n$ such that $N\in D_n$ then is it true that $f(N)=o(N)$? Perhaps just for almost all $N$?

\noindent\textbf{Related OEIS sequences:} A167485, A387502, A387503


\bigskip
\noindent\small{Source: \url{https://www.erdosproblems.com/468}}
\end{document}
